% !TEX root = ../linnik399.tex
\section{Zero location}\label{sec:location}

This section collects the lower bounds for $\lambda_1$, $\lambda'$, $\lambda_2$ and $\lambda_3$ that the
case tree uses. Most are printed results of Heath-Brown \cite{HeathBrown1992zero} and Xylouris
\cite{Xylouris2011nullstellen}; we state each in the form used, with its scope. The others are new
consequences of the same explicit formula, proved here by the positivity method; one of them uses
a sharpened conductor coefficient, which we prove in \cref{subsec:conductor}. None of the results of
this section depends on the exponent $L$.

\subsection{Conventions}\label{subsec:locconventions}
We fix once and for all an admissible height function $l(q)$ as in \cref{inp:rectangle}, for
instance the least one. All the zero-location results we quote are proved for an arbitrary
integer $l$ with $1\le l\le\LL/10$ for which $R(10l)\setminus R(l)$ contains no zero; their proofs use
$l$ only through these two properties, and their thresholds depend only on fixed data. So they
hold simultaneously for our $l(q)$.

The zeros $\rho_1,\rho_2,\rho_3$ are selected as in \cite[(3.11)--(3.12), p.~21]{Xylouris2011nullstellen}
(which follows \cite[\S6]{HeathBrown1992zero}). At step $k$, remove the characters in the
families $\mathcal F_1,\ldots,\mathcal F_{k-1}$ and choose a zero of maximal real part
in $R(l)$ among the remaining nonprincipal $L$-functions. Each real character is
removed once. This is the successive-minimum convention of \cref{def:families}.
The additional zero
$\rho'$ of the first family is selected as in \cite[pp.~21--22]{Xylouris2011nullstellen}: $\rho'=\rho_1$ if
$\rho_1$ is multiple, and otherwise $\rho'$ has maximal real part among the zeros of $L(s,\chi_1)$ in
$R(l)$ other than $\rho_1$ (and other than $\overline{\rho_1}$ if $\chi_1$ is real). Its parameter is the
$\lambda'$ of \cref{sec:notation} (Heath-Brown also writes $\lambda'$; Xylouris indexes this zero by $0$). When several zeros qualify
(ties in the real part, or several characters attaining a minimum), any of them may be selected:
the proofs of the quoted results use only the extremal property of the selection, so they hold for
every such choice. This is why the configuration sets of \cref{def:specification} quantify over the
admissible choices. All the results of this section concern zeros in $R(l)$ at any height up to
$l$, so they apply whether $\rho_1$ is inside or outside the buffer.

\emph{Printed tables.} Every printed bound is used as an exact implication for $q\ge q_0$. Each
of them was obtained by showing that a certain inequality fails with a positive margin, which
persists when the error term $\eps$ is small enough \cite[p.~43, p.~44 and footnote~2 on
p.~45]{Xylouris2011nullstellen}; for Heath-Brown's Tables~2 and~3 the printed values lie a little below
the computed roots \cite[\S8]{HeathBrown1992zero}. Heath-Brown states no convention for his
Tables~4 and~7, and we have recomputed their roots from his printed parameters; every recomputed
root exceeds the printed value, and each printed value is the root rounded down. The closest case
in Table~4 is row $1.05$, with root $1.43904$ against $1.439$, and the closest in Table~7 are rows
$0.20$ and $0.55$, with $2.01015$ against $2.01$ and $1.00036$ against $1.00$. The
rows of Table~4 are derived under the side conditions \cite[(8.6), (8.8)]{HeathBrown1992zero}; we have
checked from the printed parameters that these hold for every row from $0.35$ on (the tightest is
(8.6), with slacks $3.8\cdot10^{-6}$, $3.8\cdot10^{-6}$ and $2.3\cdot10^{-6}$ for the rows $1.05$, $1.15$ and
$1.294$). For the row $(0.3,2.293)$ they need not hold, but that row also follows from Heath-Brown's
Table~3, which gives $\lambda'\ge2.43$ for $\lambda_1\le0.30$ \cite[\S8]{HeathBrown1992zero}. These checks, and the
recomputation of the roots, were made in floating point, separately from the replay. A separate
program compares every printed number that the proof uses, in the code and in this paper, with
the text of the source pages and with the printed scope of its row; each agrees with the printed
value or is weaker than it in the safe direction.
Results stated in the form $c-\eps$ are used with an explicit reduction: $0.857<\frac67$ in
\cref{inp:lambda3}, and $1.09<\frac{12}{11}$ and $\eps=0.001$ in \cref{sec:exteriors}.

\subsection{Types and the first zero}

\begin{proposition}[Lower bounds for a nonreal first zero]\label{prop:firstlower}
For all sufficiently large $q$, the first family satisfies
\[
  \lambda_1>0.440\quad\text{if }\chi_1\text{ is nonreal},\qquad
  \lambda_1>0.628\quad\text{if }\chi_1\text{ is real and }\rho_1\text{ is nonreal}.
\]
Thus $\lambda_1\le0.44$ forces both $\chi_1$ and $\rho_1$ to be real.
These bounds are \cite[Lemma~4.5 and Table~11, pp.~60--61]{Xylouris2011nullstellen}.
\end{proposition}

Xylouris's lemma gives $\lambda_1>0.440$, $0.493$, $0.478$, $0.498$ and $0.628$ for characters of order
$\ge6$, $=5$, $=4$, $=3$ and $=2$ respectively; the minimum over orders at least $3$ is $0.440$.

\subsection{The test function}
The location inequalities use one family of test functions. We define it here so that
all parameters in the statements below are specified before they occur.
For $\gamma>0$, put
\begin{equation}\label{eq:ftest}
  f_\gamma(t)=\begin{cases}\frac{16\gamma^5}{15}(1-u)^3(1+3u+u^2),&0\le t\le2\gamma,\ u=t/(2\gamma),\\
  0,&t\ge2\gamma,\end{cases}
\end{equation}
The function $f_\gamma$ is the autocorrelation of $x\mapsto(\gamma^2-x^2)_+$, the test function of
\cite[(3.14), (3.17)--(3.18)]{Xylouris2011nullstellen}; it satisfies Conditions~1 and~2
(\cref{lem:parabolic}), and its transform $F$ is positive and decreasing on $\RR$.

\subsection{Printed zero-location inputs}

\begin{inp}[Real first zero]\label{inp:rrtables}
Let $\chi_1$ and $\rho_1$ be real. For $q\ge q_0$:
\begin{enumerate}[(a)]
  \item \cite[Tables~3 and~4, Lemmas~8.2 and~8.3]{HeathBrown1992zero}. If $\lambda_1\le t$ then
    $\lambda'\ge c$, for the $24$ pairs $(t,c)$ from $(0.3,\,2.293)$ to $(1.294,\,1.294)$ of the table.
    For a real $\rho'$ this follows from $\lambda'\ge2.427$ (Lemma~8.2).
  \item \cite[Lemma~8.4]{HeathBrown1992zero}. For every $\eps>0$: $\lambda'\ge(2-\eps)\log\lambda_1^{-1}$ if
    $\lambda_1\le0.2$ and $q\ge q(\eps)$; and $\lambda'\ge\max\bigl(\frac32\log\lambda_1^{-1},1.294\bigr)$ for every $\lambda_1$.
  \item \cite[Table~7 and Lemma~8.7]{HeathBrown1992zero}. If $\lambda_1\le t$ then $\lambda_2\ge c$, for
    the $16$ pairs $(t,c)$ from $(0.12,\,2.56)$ to $(0.745,\,0.745)$ of the table; these hold whether or
    not $\chi_2^4=\chi_0$ (the row $t=0.10$ is not used).
  \item \cite[Lemma~8.8]{HeathBrown1992zero}. $\lambda_2\ge\max\bigl((\frac{12}{11}-\eps)\log\lambda_1^{-1},0.745\bigr)$
    for $q\ge q(\eps)$.
\end{enumerate}
\end{inp}

Heath-Brown remarks that the estimates of his \S8 have not been proved for extremely small values
of $\lambda_1$ \cite[\S8]{HeathBrown1992zero}. We use (a)--(d) only for $\lambda_1$ bounded below by a
fixed positive constant: $\lambda_1\ge0.1$ in the case tree, and, in \cref{thm:small}, $\lambda_1$ at
least the fixed constant $u_{\mathrm{exc}}>0$ of \cref{sec:exteriors}.

\begin{inp}[Nonreal first character or zero]\label{inp:complextables}
Let $\chi_1$ or $\rho_1$ be nonreal. For $q\ge q_0$:
\begin{enumerate}[(a)]
  \item \cite[Table~2$'$, p.~62]{Xylouris2011nullstellen}. If $\lambda_1\le t$ then $\lambda'>c$, for the pairs
    $(0.46,1.85),\dots,(0.827,0.827)$ of the table.
  \item \cite[Table~3, p.~45]{Xylouris2011nullstellen}. If $\ord\chi_1\in\{2,3,4\}$ and $\lambda_1\le t$, then
    $\lambda'>c$, for the pairs $(0.38,2.53),\dots,(1.099,1.099)$; in particular $\lambda_1\le0.86$ implies
    $\lambda'>1.35$.
  \item \cite[Tables~6 and~7, p.~53]{Xylouris2011nullstellen}. If $\lambda_1\le t$ then $\lambda_2>c$; Table~6
    concerns a real $\chi_1$ (Fall~7), Table~7 all cases.
  \item \cite[Lemma~9.4 and Table~10]{HeathBrown1992zero}. $\lambda_2\ge0.702$ always, and
    $\lambda_1\le0.70$ implies $\lambda_2\ge0.704$.
\end{enumerate}
\end{inp}

\begin{inp}[The third family]\label{inp:lambda3}
For $q\ge q_0$:
\begin{enumerate}[(a)]
  \item \cite[Lemma~10.3]{HeathBrown1992zero}. $\lambda_3\ge\frac67-\eps$; we use $\lambda_3>0.857$.
  \item \cite[Table~8, column ``alle F\"alle'', p.~55]{Xylouris2011nullstellen}. If $\chi_1$ or $\rho_1$ is
    nonreal and $\lambda_1\le t$, then $\lambda_3>c$, for $(t,c)=(0.52,1.320)$, $(0.54,1.243)$,
    $(0.56,1.160)$, $(0.58,1.079)$, $(0.60,1.001)$, $(0.62,0.933)$.
  \item \cite[Lemma~4.4 and Table~10, p.~55]{Xylouris2011nullstellen}. If $\chi_1$ and $\rho_1$ are real, then
    $\lambda_1\in[0.44,0.60]$, $[0.60,0.70]$, $[0.70,0.80]$ imply $\lambda_3>1.176$, $1.055$, $0.952$.
  \item \cite[(4.28)--(4.29), pp.~56--57]{Xylouris2011nullstellen}. Let $\chi_1$ be nonreal and
    $\lambda_1\in[0.44,0.85]$, and let $f$ be the test function \eqref{eq:ftest} with $\gamma=\frac54$. Then
    \[
      0\le F(-\lambda_1)-F(\lambda_3-\lambda_1)-F(\lambda_2-\lambda_1)-F(0)+\tfrac76f(0)+\eps,
    \]
    provided that $\sup_{y\in\RR}\Re\{F(-\lambda_1+iy)-2F(iy)\}\le\frac16f(0)$.
  \item \cite[(4.31) and (4.34), pp.~58--59]{Xylouris2011nullstellen}. Let $\chi_1$ and $\rho_1$ be real,
    $\lambda_1\in[0.44,0.80]$ and $\lambda_2\in[0.44,1.176]$, and let $f$ be \eqref{eq:ftest} with
    $\gamma=1.04$. Then
    \[
      0\le F(-\lambda_2)-F(\lambda_3-\lambda_2)-F(0)-F(\lambda_1-\lambda_2)+\tfrac98f(0)+\eps,
    \]
    provided that $\sup_{t\in\RR}\Re\{F(-\lambda_2+it)-F(\lambda_1-\lambda_2+it)-F(it)\}\le\frac5{48}f(0)$ on this box.
\end{enumerate}
\end{inp}

In (d), the side condition is \cite[(4.29)]{Xylouris2011nullstellen}, needed when a product character in his argument
is principal; he verified it numerically, and we have verified it with interval arithmetic for
all $\lambda_1\in[0.44,0.85]$: the supremum is at most $0.00335$, while $\frac16f(0)=0.5425\ldots$. In (e),
Xylouris's argument \cite[(4.32)--(4.34), p.~59]{Xylouris2011nullstellen} needs
\[
  \sup_{t\in\RR}\Re\bigl\{F(-\lambda_2+it)-F(\lambda_1-\lambda_2+it)-F(it)\bigr\}\le\tfrac5{48}f(0)
\]
on the whole box, and he states that the supremum is below $0.10$, while $\frac5{48}f(0)=0.1351\ldots$. We have
verified this with interval arithmetic as well: over $\lambda_1\in[0.44,0.80]$, $\lambda_2\in[0.44,1.18]$ and all real
$t$ the supremum is at most $0.0031$. (Both checks enclose the transforms on a grid of heights with
step $\frac1{100}$ up to $30$, with second-derivative interpolation errors in all variables, and use a
closed-form bound beyond height $30$.) (Xylouris derives (4.28) from Heath-Brown's (10.2)--(10.6), whose derivation does not use
the assumption $\lambda_3\le\frac67$ made in \cite[\S10]{HeathBrown1992zero}; he himself applies it to
prove bounds above $\frac67$.)

\begin{proposition}[Parent rows]\label{prop:parentrows}
The $58$ \emph{parent rows} $(\mathfrak t,A_j,B_j,p_j,r_j)$ of \cref{tab:parents}, from which the roots of
the case tree are built (\cref{subsec:sourcecover}), are valid implications: for $q\ge q_0$, every
configuration of type $\mathfrak t$ with $\lambda_1\in[A_j,B_j]$ has $\lambda'\ge p_j$ and $\lambda_2\ge r_j$.
\end{proposition}

\begin{proof}
Each entry is at most the largest of the following bounds, each valid on the whole interval $[A_j,B_j]$ for
the type $\mathfrak t$: (i) the applicable rows $(t,c)$ of \cref{inp:rrtables,inp:complextables} with $t\ge B_j$; (ii) the
bound $c$ of an applicable row $(t,c)$ with $c=t$ (for $\lambda_1\le t$ the row gives it, and for
$\lambda_1>t$ the trivial bound gives $\lambda',\lambda_2\ge\lambda_1>t=c$), as for the rows
$(0.827,0.827)$ of Table~$2'$ and $(1.099,1.099)$ of Table~3 of \cite{Xylouris2011nullstellen} and $(0.745,0.745)$
of Table~7 of \cite{HeathBrown1992zero}; (iii) the bounds of \cref{inp:rrtables}(b),(d) and
\cref{inp:complextables}(d), which hold for every $\lambda_1$; and (iv) the trivial bounds
$\lambda',\lambda_2\ge\lambda_1\ge A_j$.
\end{proof}

\begin{table}[tp]
\centering
\small
\begin{tabular}{@{}lccc@{\qquad\qquad}lccc@{}}
\toprule
Type & $[A_j,B_j]$ & $p_j$ & $r_j$ & Type & $[A_j,B_j]$ & $p_j$ & $r_j$\\
\midrule
$\mathrm{rr}$ & $[0.1,0.12]$ & $2.293$ & $2.56$ & $\mathrm{rc}$ & $[0.628,0.66]$ & $1.67$ & $0.93$ \\
$\mathrm{rr}$ & $[0.12,0.14]$ & $2.293$ & $2.39$ & $\mathrm{rc}$ & $[0.66,0.7]$ & $1.59$ & $0.93$ \\
$\mathrm{rr}$ & $[0.14,0.16]$ & $2.293$ & $2.25$ & $\mathrm{rc}$ & $[0.7,0.74]$ & $1.52$ & $0.93$ \\
$\mathrm{rr}$ & $[0.16,0.18]$ & $2.293$ & $2.12$ & $\mathrm{rc}$ & $[0.74,0.78]$ & $1.46$ & $0.93$ \\
$\mathrm{rr}$ & $[0.18,0.2]$ & $2.293$ & $2.01$ & $\mathrm{rc}$ & $[0.78,0.82]$ & $1.40$ & $0.82$ \\
$\mathrm{rr}$ & $[0.2,0.25]$ & $2.293$ & $1.77$ & $\mathrm{rc}$ & $[0.82,0.86]$ & $1.35$ & $0.82$ \\
$\mathrm{rr}$ & $[0.25,0.3]$ & $2.293$ & $1.58$ & $\mathrm{rc}$ & $[0.86,0.9]$ & $1.30$ & $0.82$ \\
$\mathrm{rr}$ & $[0.3,0.348]$ & $2.195$ & $1.42$ & $\mathrm{rc}$ & $[0.9,0.94]$ & $1.25$ & $0.82$ \\
$\mathrm{rr}$ & $[0.348,0.35]$ & $2.195$ & $1.42$ & $\mathrm{rc}$ & $[0.94,0.98]$ & $1.21$ & $0.82$ \\
$\mathrm{rr}$ & $[0.35,0.4]$ & $2.108$ & $1.29$ & $\mathrm{rc}$ & $[0.98,1.02]$ & $1.17$ & $0.82$ \\
$\mathrm{rr}$ & $[0.4,0.45]$ & $2.03$ & $1.18$ & $\mathrm{rc}$ & $[1.02,1.06]$ & $1.13$ & $0.82$ \\
$\mathrm{rr}$ & $[0.45,0.5]$ & $1.958$ & $1.08$ & $\mathrm{rc}$ & $[1.06,1.099]$ & $1.099$ & $0.82$ \\
$\mathrm{rr}$ & $[0.5,0.55]$ & $1.893$ & $1.0$ & $\mathrm{rc}$ & $[1.099,1.5]$ & $1.099$ & $0.82$ \\
$\mathrm{rr}$ & $[0.55,0.6]$ & $1.832$ & $0.92$ & $\mathrm{complex}$ & $[0.44,0.58]$ & $1.36$ & $1.04$ \\
$\mathrm{rr}$ & $[0.6,0.65]$ & $1.776$ & $0.85$ & $\mathrm{complex}$ & $[0.58,0.64]$ & $1.15$ & $0.85$ \\
$\mathrm{rr}$ & $[0.65,0.7]$ & $1.724$ & $0.79$ & $\mathrm{complex}$ & $[0.64,0.66]$ & $1.08$ & $0.79$ \\
$\mathrm{rr}$ & $[0.7,0.745]$ & $1.676$ & $0.745$ & $\mathrm{complex}$ & $[0.66,0.68]$ & $1.02$ & $0.74$ \\
$\mathrm{rr}$ & $[0.745,0.75]$ & $1.676$ & $0.745$ & $\mathrm{complex}$ & $[0.68,0.70]$ & $0.96$ & $0.704$ \\
$\mathrm{rr}$ & $[0.75,0.8]$ & $1.63$ & $0.745$ & $\mathrm{complex}$ & $[0.70,0.72]$ & $0.93$ & $0.702$ \\
$\mathrm{rr}$ & $[0.8,0.85]$ & $1.587$ & $0.8$ & $\mathrm{complex}$ & $[0.72,0.74]$ & $0.91$ & $0.72$ \\
$\mathrm{rr}$ & $[0.85,0.9]$ & $1.547$ & $0.8$ & $\mathrm{complex}$ & $[0.74,0.76]$ & $0.89$ & $0.74$ \\
$\mathrm{rr}$ & $[0.9,0.95]$ & $1.509$ & $0.8$ & $\mathrm{complex}$ & $[0.76,0.78]$ & $0.86$ & $0.74$ \\
$\mathrm{rr}$ & $[0.95,1]$ & $1.473$ & $0.8$ & $\mathrm{complex}$ & $[0.78,0.80]$ & $0.84$ & $0.78$ \\
$\mathrm{rr}$ & $[1,1.05]$ & $1.439$ & $0.8$ & $\mathrm{complex}$ & $[0.80,0.82]$ & $0.83$ & $0.78$ \\
$\mathrm{rr}$ & $[1.05,1.1]$ & $1.406$ & $0.8$ & $\mathrm{complex}$ & $[0.82,1.5]$ & $0.827$ & $0.82$ \\
$\mathrm{rr}$ & $[1.1,1.15]$ & $1.375$ & $0.8$ &  & & &  \\
$\mathrm{rr}$ & $[1.15,1.175]$ & $1.36$ & $0.8$ &  & & &  \\
$\mathrm{rr}$ & $[1.175,1.2]$ & $1.346$ & $0.8$ &  & & &  \\
$\mathrm{rr}$ & $[1.2,1.225]$ & $1.331$ & $0.8$ &  & & &  \\
$\mathrm{rr}$ & $[1.225,1.25]$ & $1.318$ & $0.8$ &  & & &  \\
$\mathrm{rr}$ & $[1.25,1.275]$ & $1.304$ & $0.8$ &  & & &  \\
$\mathrm{rr}$ & $[1.275,1.294]$ & $1.294$ & $0.8$ &  & & &  \\
$\mathrm{rr}$ & $[1.294,1.5]$ & $1.294$ & $0.8$ &  & & &  \\
\bottomrule
\end{tabular}
\caption{Parent implications: for the stated type and $\lambda_1\in[A_j,B_j]$,
$\lambda'\ge p_j$ and $\lambda_2\ge r_j$. Each panel is read downwards, and each row is
valid on the whole closed interval. The left panel gives the $33$ real-first-zero rows;
the right gives the $13$ real-character/nonreal-zero rows and $12$ nonreal-character rows.}
\label{tab:parents}
\end{table}

\subsection{The positivity method}
The use of nonnegative trigonometric polynomials to locate zeros goes back to de la Vall\'ee
Poussin \cite{delaValleePoussin1899fonction}; it underlies all zero-free regions for Dirichlet
$L$-functions \cite{Gronwall1913series,Landau1918imaginar,Page1935number,McCurley1984explicit,
Kadiri2018explicit} and for $\zeta$ \cite{Kadiri2005region,MossinghoffTrudgian2015nonnegative}, and,
combined with smoothed explicit formulas, the tables of Heath-Brown and Xylouris and their
refinements \cite{LiuWang1998numerical}. All the new zero-location results come from the following
consequence of
\cite[Lemmas~3.1 and~3.4]{Xylouris2011nullstellen}. Lemma~3.4 there is a ``working version'' of
\cref{inp:X32}: if $f$ satisfies Conditions~1 and~2, $\chi\ne\chi_0$ and $s\in R(9l)$, and $A_1$ is the
set of zeros $\rho\in R(l)$ of $L(s,\chi)$ with $\Re\rho>\Re s$ (with multiplicity, at most $N$ of them)
and $A_2$ is any set of at most $N$ zeros $\rho\in R(l)$ with $\Re\rho\le\Re s$, then for every $\eps>0$
and $q\ge q_0(f,\eps,N)$
\begin{equation}\label{eq:X34}
  \sum_n\Lambda(n)\Re\frac{\chi(n)}{n^s}f\Bigl(\frac{\log n}\LL\Bigr)\le
  -\LL\sum_{\rho\in A_1\cup A_2}\Re F\bigl((s-\rho)\LL\bigr)+\frac{\varphi(\chi)}2f(0)\LL+\eps\LL .
\end{equation}

\begin{lemma}[Positivity method]\label{lem:positivity}
Let $f$ satisfy Conditions~1 and~2, with Laplace transform $F$. Fix a nonnegative
anchor $a$, and suppose $a\le\lambda_1$. Let finitely many triples
$(c_j,\chi^{(j)},t_j)$ be given, with $c_j>0$, $\chi^{(j)}$ a character modulo $q$ and $|t_j|\le9l$, such
that
\[
  \sum_jc_j\Re\bigl(\chi^{(j)}(n)n^{-it_j}\bigr)\ge0\qquad\text{for all $n$ coprime to $q$.}
\]
For each nonprincipal $\chi^{(j)}$ let $A_2(j)$ be a set of at most $N$ zeros of $L(s,\chi^{(j)})$ in $R(l)$,
with multiplicity. Then for every $\eps>0$ and $q\ge q_0$,
\[
  \sum_{\chi^{(j)}=\chi_0}c_j\Re F(-a+it_j\LL)-\sum_{\chi^{(j)}\ne\chi_0}c_j\sum_{\rho\in A_2(j)}\Re
  F\bigl((1+it_j-\rho)\LL-a\bigr)+\frac{f(0)}2\sum_{\chi^{(j)}\ne\chi_0}c_j\varphi(\chi^{(j)})+\eps\ge0 .
\]
\end{lemma}

\begin{proof}
Multiply the hypothesis by $\Lambda(n)n^{-(1-a/\LL)}f(\log n/\LL)\ge0$ and sum over $n$ (for $(n,q)>1$
every term vanishes, the principal ones included). For principal $\chi^{(j)}$ use \cref{inp:X31} at
$1-a/\LL+it_j$. For the others use \eqref{eq:X34} at $s_j=1-a/\LL+it_j$, which lies in $R(9l)$;
there $A_1=\emptyset$, because every zero in $R(l)$ of a nonprincipal $L$-function has parameter at
least $\lambda_1\ge a$. Divide by $\LL$.
\end{proof}

In every application the retained zeros satisfy $\lambda_\rho\ge a$, the transform $F$ is decreasing on
$\RR$, and the terms that we cannot evaluate are bounded over all real heights with interval
arithmetic. A necessary inequality between the parameters of the retained zeros results, and a
case is excluded when that inequality fails with a positive margin. Since there are finitely many
such inequalities, each with a fixed positive margin, one $q_0$ serves all of them.

\subsection{A sharper conductor coefficient}\label{subsec:conductor}
Heath-Brown remarks \cite[\S16, item~4]{HeathBrown1992zero} that his estimate for $\lambda_2$ in the
case of a real first character can be improved by replacing the conductor coefficient $\frac{11}{24}$
by $\frac{97}{216}$; he describes the ingredients but gives no proof. We need a weighted form of this
refinement, and we prove it here. For a modulus $q$ put
\[
  v=\prod_{p^e\parallel q,\ e\ge3}p^e,\qquad \theta_q=\frac{\log v}{\log q}\in[0,1] .
\]
We also write $\operatorname{rad}q=\prod_{p\mid q}p$.

\begin{lemma}[Growth bounds with conductor dependence]\label{lem:Lbounds}
Let $k\ge3$ and $\eps>0$. For every nonprincipal character $\chi$ modulo $q$,
\[
  |L(\sigma+it,\chi)|\ll_{\eps,k}q^{\phi_q(\chi)(1-\sigma)(1+k^{-1})+\eps}(1+|t|)\qquad
  \Bigl(1-\frac1k\le\sigma\le1+\frac{\log\LL}\LL\Bigr),
\]
where $\phi_q(\chi)=\frac14\frac{\log(4\operatorname{rad}q)}{\log q}$ if $\chi$ is real and
$\phi_q(\chi)=\min\bigl(\frac13,\frac14+\frac34\theta_q\bigr)$ otherwise.
\end{lemma}

\begin{proof}
Let $\chi$ be induced by the primitive character $\chi^*$ modulo $q^*$; then $q^*\mid q$ and
$|L(s,\chi)|\ll q^\eps|L(s,\chi^*)|$ for $\sigma\ge\frac12$ \cite[p.~10]{HeathBrown1992zero}. Heath-Brown
proves (\cite[pp.~9--10]{HeathBrown1992zero}, (2.2)--(2.4)) that a bound
$\sum_{N<n\le N+H}\chi^*(n)\ll_{\eps,k}Q^{(k+1)/(4k^2)}q^{*\eps}H^{1-1/k}$ for $1\le H\le q^*$ gives, by partial
summation with the cut $q_1=Q^{(k+1)/(4k)}$ and the P\'olya--Vinogradov inequality beyond $q^*$,
$L(s,\chi^*)\ll Q^{(1-\sigma)(1+k^{-1})/4}q^{*\eps}(1+|t|)$ in the stated range (if $q_1>q^*$ the trivial
bound up to $q^*$ is already stronger). For $\sigma\le1$ a bound of this form with a smaller $Q$ is
stronger; for $1<\sigma\le1+\log\LL/\LL$ every factor $Q^{(1-\sigma)c}$ with $1\le Q\le q^4$ and
$0\le c\le1$ lies between $\LL^{-4}$ and $1$, so all such bounds agree up to a factor $q^{\eps}$.

If $\chi$ is real, then $\chi^*$ is a real primitive character, so $q^*$ is the absolute value of a
fundamental discriminant and $q^*\le4\operatorname{rad}(q^*)\le4\operatorname{rad}(q)$. Its order is $2$, so by \cite[Lemma~2.4]{HeathBrown1992zero}, which bounds the
sums with the factor $l^{3/(2k)}$ for a character of order $l$, the character-sum bound holds with
$Q=q^*$ for every $k$. This gives the real case.

In general, \cite[Lemma~2.3]{HeathBrown1992zero} gives the bound with
$Q^{(k+1)/(4k^2)}=v^{*(3k-1)/(4k^2)}q^{*(k+1)/(4k^2)}\le(v^{*3}q^*)^{(k+1)/(4k^2)}$, where $v^*$ is the
cube-full part of $q^*$; since $q^*\mid q$ we have $v^*\mid v$ and $v^{*3}q^*\le v^3q$. This gives the
exponent $\frac14(1+\frac{3\log v}{\log q})=\frac14+\frac34\theta_q$. The exponent $\frac13$ is
\cite[Lemma~2.5]{HeathBrown1992zero}. Taking the smaller bound gives the general case. (Bounds that are uniform in $q$ and $t$
together are studied in \cite{HeathBrown1978hybrid}; here the factor $1+|t|$ is harmless, since
$|t|\le\LL+1$.)
\end{proof}

\begin{lemma}[The refined coefficient in the explicit formula]\label{lem:explicitphi}
For every $\eps>0$, \cref{inp:X32} and \eqref{eq:X34} remain true, with a suitable $q_0$, when $\varphi(\chi)$ is replaced by
$\phi_q(\chi)+\eps$, with the radius $\delta$ that the proof of \cref{inp:X32} provides for $\varphi=\frac13$ (which
serves all characters).
\end{lemma}

\begin{proof}
Both are deduced from Heath-Brown's Lemma~3.1 \cite[pp.~12--14, 23--24]{HeathBrown1992zero}
\cite[pp.~19--21]{Xylouris2011nullstellen}, whose proof uses the bound (2.5) of his Lemma~2.5 only through
the inequality
\[
  \log|L(s_0+Re^{i\vartheta},\chi)|\le\bigl\{\phi(1+k^{-1})(1-\sigma_0-R\cos\vartheta)+2\eps\bigr\}\LL
\]
for $\frac\pi2\le\vartheta\le\frac{3\pi}2$, $R\le1/k$ and large $q$. By \cref{lem:Lbounds} this holds with
$\phi=\phi_q(\chi)$, with an implied constant independent of $q$ and $\chi$. The remaining choices in
that proof ($k=3+[3\phi/(2\eps_0)]$, $R=1/k$, and the radius) may be made with the upper bound
$\phi\le\frac13$ in place of $\phi$, so they do not depend on $q$. The deduction of Lemma~5.2 of
\cite{HeathBrown1992zero} (\cref{inp:X32}) and of \eqref{eq:X34} from Lemma~3.1 is unchanged.
\end{proof}

\begin{proposition}[A coupled saving in the conductor term]\label{prop:conductor}
Let $\chi_1$ be a real nonprincipal character and let $\chi_2$ be a character with $\chi_2^k\ne\chi_0$ and $\chi_1\chi_2^k\ne\chi_0$
for $k\in K$, a finite set of positive integers. Let $c_k>0$ ($k\in K$) with $\Sigma=\sum_kc_k\ge\frac19$.
Suppose the polynomial $1+\sum_{k\in K}c_k\cos(k\vartheta)$ is nonnegative for every real
$\vartheta$. In an application of \cref{lem:positivity} to
$(1+\chi_1(n))\bigl(1+\sum_{k\in K}c_k\Re\bigl((\chi_2(n)n^{-i\gamma_2})^k\bigr)\bigr)$,
the conductor term may, for every fixed $\eps>0$ and sufficiently large $q$, be taken to be
\[
  \Bigl(\frac18+\frac\Sigma3-\frac1{108}\Bigr)f(0)+\eps .
\]
\end{proposition}

\begin{proof}
The nonprincipal characters are $\chi_1$ with coefficient $1$, and $\chi_2^k$ and $\chi_1\chi_2^k$, each with
coefficient $c_k$. By \cref{lem:explicitphi} the conductor term is at most
$\frac{f(0)}2\bigl(\phi_q(\chi_1)+2\Sigma\max_\chi\phi_q(\chi)\bigr)+\eps$, the maximum being over the
characters $\chi_2^k$ and $\chi_1\chi_2^k$. Since $p\le(p^e)^{1/3}$ when $e\ge3$,
$\operatorname{rad}q\le qv^{-2/3}$, so $\phi_q(\chi_1)\le\frac14-\frac{\theta_q}6+o(1)$, and
$\phi_q(\chi)\le\min(\frac13,\frac14+\frac34\theta_q)+o(1)$ for every nonprincipal $\chi$ (for a real $\chi$
because $\frac14-\frac{\theta_q}6\le\frac14$).
The conductor term is therefore at most $\mathcal K(\theta_q)f(0)+2\eps$ for large $q$, with
\[
  \mathcal K(x)=\frac18-\frac x{12}+\Sigma\min\Bigl(\frac13,\frac14+\frac{3x}4\Bigr)\qquad(0\le x\le1).
\]
On $[0,\frac19]$ the slope of $\mathcal K$ is $\frac34\Sigma-\frac1{12}\ge0$, and on $[\frac19,1]$ it is
$-\frac1{12}$. Hence $\mathcal K(\theta_q)\le\mathcal K(\frac19)=\frac18+\frac\Sigma3-\frac1{108}$.
\end{proof}

For $\Sigma=1$ this is Heath-Brown's $\frac{97}{216}=\frac{11}{24}-\frac1{108}$.

\subsection{Lower bounds for \texorpdfstring{$\lambda_2$}{lambda2} when the first zero is real}

\begin{proposition}[Second-family bounds for a real first zero]\label{prop:realrows}
Let $\chi_1,\rho_1$ be real with $\lambda_1\in[a,b]$, where $0\le a\le b$, and let $h\ge b$. Let
$P(x)=1+c_1x+c_2(2x^2-1)+c_5(16x^5-20x^3+5x)$, that is
$P(\cos\vartheta)=1+c_1\cos\vartheta+c_2\cos2\vartheta+c_5\cos5\vartheta$, with
\[
  (c_1,c_2,c_5)=\frac{(1.41866466,\ 0.43287503,\ 0.01421037)}{1.000001},
\]
so that $P\ge0$ on $[-1,1]$ and $\Sigma=c_1+c_2+c_5\ge\frac19$.
Choose $f=f_\gamma$ from \eqref{eq:ftest}, with transform $F$ and fixed $\gamma>0$.
If $\lambda_2\le h$, then for every $\eps>0$ and all sufficiently large $q$, the inequality
corresponding to the character relation below must hold. The generic case means that
$\chi_2$ is nonreal and none of the relations in (ii)--(iii) holds.
\begin{enumerate}[(i)]
  \item (generic) $F(b-a)+c_1F(h-a)\le F(-a)+(\frac18+\frac\Sigma3-\frac1{108})f(0)+\eps$;
  \item ($\chi_2^2=\chi_1$) $F(b-a)+c_1F(h-a)\le F(-a)+(\frac{1+c_2}8+\frac{c_1+c_5}3)f(0)+\sup_\Delta R(\Delta)+\eps$,
    where $R(\Delta)=c_2\{\Re F(-a+i\Delta)-\Re F(\lambda_1-a+i\Delta)\}-c_1\Re F(\lambda_2-a+i\Delta)$, the supremum
    over $\Delta\in\RR$, $\lambda_1\in[a,b]$ and $\lambda_2\in[a,h]$;
  \item ($\chi_2^5\in\{\chi_0,\chi_1\}$) $F(b-a)+c_1F(h-a)\le(1+c_5)F(-a)+(\frac{1+c_5}8+\frac{c_1+c_2}4)f(0)+\eps$;
  \item ($\chi_2$ real) $F(b-a)+F(h-a)\le F(-a)+\frac38f(0)+\eps$, with a second test function
    \eqref{eq:ftest}.
\end{enumerate}
Here $\chi_2$ is a character of the second family. The test parameter in (iv) may be
chosen independently of the one in (i)--(iii). If all four possible cases are excluded
with a positive margin at $\eps=0$, then $\lambda_2>h$ for all sufficiently large $q$.
\end{proposition}

\begin{proof}
Put $z_n=\chi_2(n)n^{-i\gamma_2}$ and apply \cref{lem:positivity} with anchor $a$ to
$(1+\chi_1(n))P(\Re z_n)\ge0$ (with $\Re z_n^k=T_k(\Re z_n)$ for $|z_n|=1$), retaining $\rho_1$ in the term of $\chi_1$
and $\rho_2$ in the term $c_1\Re z_n$. The heights are at most $5l\le9l$. The nonzero frequencies are
$1$, $2$ and $5$, and $\chi_2\ne\chi_1$. A product character is principal only if $\chi_2$ is real, or
$\chi_2^2=\chi_1$ (order $4$), or $\chi_2^5\in\{\chi_0,\chi_1\}$ (order $5$ or $10$); the last two cannot
occur together, since together they would give $\chi_2^2=\chi_0$. Generic case: there is no principal
product, and \cref{prop:conductor} gives (i), using that $F$ is decreasing, $\lambda_1\le b$ and $\lambda_2\le h$.
Case $\chi_2^2=\chi_1$: then $\chi_1\chi_2^2=\chi_0$ at height $2\gamma_2$, $\chi_2^2=\chi_1$ at height $2\gamma_2$
retains $\rho_1$, and $\chi_1\chi_2=\bchi_2$ retains $\overline{\rho_2}$; with $\Delta=2\mu_2$ these give $R(\Delta)$,
while the characters of frequencies $1$ and $5$ are nonprincipal and are charged $\frac13$ (conservatively,
since some of them have order $4$). Case
$\chi_2^5\in\{\chi_0,\chi_1\}$: exactly one term is principal, bounded by $c_5F(-a)$ since $f\ge0$, and
every other character has order at most $10\le\LL$, so $\varphi=\frac14$. Real $\chi_2$: apply the
method to $(1+\chi_1(n))(1+\Re z_n)\ge0$, whose three nonprincipal characters are real.
\end{proof}

The rows used are given in \cref{tab:realrows}: $22$ cells of width $0.0025$ covering $[0.700,0.755]$.
For each cell $[a,b]$ we give $h$ and the test parameter $\gamma$ of the generic case; each of the four
inequalities fails by the margin shown (normalized by $f(0)$), so that $\lambda_2>h$ throughout the
cell. The margins were computed with interval arithmetic: the correlation term $R(\Delta)$ is enclosed
on a grid in $\Delta\in[0,30]$ with second-derivative interpolation errors in all variables, and by a
closed-form bound for $|\Delta|\ge30$. The nonnegativity of $P$ was verified at the points
$x=j/2000$ with a second-derivative bound, and independently by exact Bernstein expansion. In
addition, on cells inside $[0.700,0.7025]$ we use the degree-$2$ row
$P(\cos\vartheta)=(\frac78+\cos\vartheta)^2/\frac{81}{64}$ without the conductor refinement, which gives
$\lambda_2>0.762$ (\cref{prop:realfirstsecond}).

\begin{proposition}[A second-family bound near $0.70$]\label{prop:realfirstsecond}
If $\chi_1,\rho_1$ are real and $0.700\le\lambda_1\le0.7025$, then $\lambda_2>0.762$ for $q\ge q_0$.
\end{proposition}

\begin{proof}
As in \cref{prop:realrows} with $c_1=\frac{112}{81}$, $c_2=\frac{32}{81}$, $c_5=0$, the printed conductor
coefficients and $\gamma=1.09$ (and $\gamma=0.82$ for a real $\chi_2$); the normalized margins of the
generic, order-$4$ and real cases are at least $3.6\cdot10^{-4}$, $0.082$ and $0.033$.
\end{proof}

\begin{table}[ht]
\centering
\scriptsize
\begin{tabular}{@{}cccccccc@{}}
\toprule
cell $[a,b]$ & $h$ & $\gamma$ & generic & order $4$ & order $5,10$ & real & \\
\midrule
$[.7000,.7025]$ & $.7793$ & $1.09624$ & $1.02\cdot10^{-4}$ & $.0804$ & $.1280$ & $.0289$\\
$[.7025,.7050]$ & $.7783$ & $1.09522$ & $5.86\cdot10^{-5}$ & $.0803$ & $.1279$ & $.0287$\\
$[.7050,.7075]$ & $.7772$ & $1.09419$ & $8.66\cdot10^{-5}$ & $.0804$ & $.1280$ & $.0285$\\
$[.7075,.7100]$ & $.7762$ & $1.09317$ & $4.92\cdot10^{-5}$ & $.0803$ & $.1279$ & $.0284$\\
$[.7100,.7125]$ & $.7751$ & $1.09216$ & $8.35\cdot10^{-5}$ & $.0804$ & $.1280$ & $.0282$\\
$[.7125,.7150]$ & $.7741$ & $1.09115$ & $5.21\cdot10^{-5}$ & $.0803$ & $.1279$ & $.0280$\\
$[.7150,.7175]$ & $.7730$ & $1.09014$ & $9.25\cdot10^{-5}$ & $.0804$ & $.1279$ & $.0278$\\
$[.7175,.7200]$ & $.7720$ & $1.08914$ & $6.71\cdot10^{-5}$ & $.0804$ & $.1279$ & $.0276$\\
$[.7200,.7225]$ & $.7710$ & $1.08814$ & $4.46\cdot10^{-5}$ & $.0803$ & $.1279$ & $.0275$\\
$[.7225,.7250]$ & $.7699$ & $1.08715$ & $9.41\cdot10^{-5}$ & $.0804$ & $.1279$ & $.0273$\\
$[.7250,.7275]$ & $.7689$ & $1.08615$ & $7.75\cdot10^{-5}$ & $.0804$ & $.1279$ & $.0271$\\
$[.7275,.7300]$ & $.7679$ & $1.08517$ & $6.39\cdot10^{-5}$ & $.0804$ & $.1279$ & $.0269$\\
$[.7300,.7325]$ & $.7669$ & $1.08418$ & $5.31\cdot10^{-5}$ & $.0804$ & $.1279$ & $.0267$\\
$[.7325,.7350]$ & $.7659$ & $1.08321$ & $4.53\cdot10^{-5}$ & $.0804$ & $.1279$ & $.0265$\\
$[.7350,.7375]$ & $.7649$ & $1.08223$ & $4.03\cdot10^{-5}$ & $.0804$ & $.1278$ & $.0263$\\
$[.7375,.7400]$ & $.7639$ & $1.08126$ & $3.83\cdot10^{-5}$ & $.0804$ & $.1278$ & $.0261$\\
$[.7400,.7425]$ & $.7629$ & $1.08029$ & $3.91\cdot10^{-5}$ & $.0804$ & $.1278$ & $.0260$\\
$[.7425,.7450]$ & $.7619$ & $1.07933$ & $4.28\cdot10^{-5}$ & $.0804$ & $.1278$ & $.0258$\\
$[.7450,.7475]$ & $.7609$ & $1.07837$ & $4.93\cdot10^{-5}$ & $.0804$ & $.1278$ & $.0256$\\
$[.7475,.7500]$ & $.7599$ & $1.07741$ & $5.87\cdot10^{-5}$ & $.0804$ & $.1278$ & $.0254$\\
$[.7500,.7525]$ & $.7589$ & $1.07645$ & $7.09\cdot10^{-5}$ & $.0804$ & $.1278$ & $.0252$\\
$[.7525,.7550]$ & $.7579$ & $1.07550$ & $8.60\cdot10^{-5}$ & $.0804$ & $.1278$ & $.0250$\\
\bottomrule
\end{tabular}
\caption{The degree-$5$ rows of \cref{prop:realrows}: on each cell, $\lambda_2>h$. The values of $\gamma$
are rounded to five decimals (the exact rational values are in the data). The margins are
normalized by $f(0)$. The real case uses $\gamma=0.82$.}
\label{tab:realrows}
\end{table}

\subsection{Lower bounds for \texorpdfstring{$\lambda_2$}{lambda2} and \texorpdfstring{$\lambda'$}{lambda'} when the first character is nonreal}
For $t>0$ put
\[
  P_t(e^{i\vartheta})=\frac{(t+\cos\vartheta)^2}{t^2+\frac12}=1+c_1\cos\vartheta+c_2\cos2\vartheta,\qquad
  c_1=\frac{2t}{t^2+\frac12},\quad c_2=\frac{1/2}{t^2+\frac12},
\]
so that $c_1/c_2=4t$.

\begin{proposition}[Second-family bounds for a nonreal first character]\label{prop:complexrows}
Let $\chi_1$ be nonreal, let $\lambda_1\in[a,b]$ with $0\le a\le b$, and let $h\ge b$.
Fix $t>0$, let $c_1,c_2$ be the coefficients of $P_t$ above, and let
$f=f_\gamma$ with $\gamma>0$ and transform $F$.
Suppose $\lambda_2\le h$. For every $\eps>0$ and sufficiently large $q$, the following
necessary inequalities hold:
\begin{enumerate}[(i)]
  \item if $\chi_2$ is nonreal and
    \begin{equation}\label{eq:alias}
      \sup_{\lambda\in[a,h],\,y\in\RR}\Re\{F(-a+iy)-4tF(\lambda-a+iy)\}\le\tfrac16f(0),
    \end{equation}
    then $c_1\{F(b-a)+F(h-a)\}\le F(-a)+\frac{(1+c_1+c_2)^2-1}6f(0)+\eps$;
  \item if $\chi_2$ is real, choose $t'>0$ and let $c_1',c_2'$ be the coefficients of $P_{t'}$.
    If \eqref{eq:alias} holds with $t'$ in place of $t$ and $[a,b]$ in place of
    $[a,h]$, then $c_1'F(b-a)+F(h-a)\le F(-a)+(\frac18+\frac{c_1'+c_2'}3)f(0)+\eps$.
\end{enumerate}
\end{proposition}

\begin{proof}
(i) Apply \cref{lem:positivity} to $P_t(\chi_1(n)n^{-i\gamma_1})P_t(\chi_2(n)n^{-i\gamma_2})\ge0$, expanded as
\[
  \sum_{w\in\{-2,\dots,2\}^2}a_{|w_1|}a_{|w_2|}\,\chi^w(n)\,n^{-iw\cdot\gamma},\qquad
  a_0=1,\ a_1=\tfrac{c_1}2,\ a_2=\tfrac{c_2}2,\ \chi^w=\chi_1^{w_1}\chi_2^{w_2}.
\] The indices $(\pm1,0)$ retain $\rho_1$ (or $\overline{\rho_1}$) and $(0,\pm1)$ retain $\rho_2$
(or $\overline{\rho_2}$). If no nonzero index is principal, the conductor term is at most
$\frac16((1+c_1+c_2)^2-1)f(0)$. A nonzero principal index $w$ is not on an axis, since $\chi_1$ and $\chi_2$
are nonreal, and does not have $|w_1|,|w_2|\le1$, since $\chi_2\notin\{\chi_1,\bchi_1\}$. So some $|w_i|=2$;
choose such $i$, with $i=1$ in case of a tie, and put $T(w)=w-\operatorname{sgn}(w_i)e_i$. Then
$\chi^{T(w)}=\chi_i^{-\operatorname{sgn}w_i}$ is nonprincipal and not an axis index of length one, its
coefficient is $4t$ times that of $w$, and the zero of $\chi_i^{-\operatorname{sgn}w_i}$ selected at the axis
index sits at relative height $w\cdot\mu$. The map $T$ is injective: a collision could only come from
$(2s,\sigma)$ and $(s,2\sigma)$ with $s,\sigma=\pm1$, and if both were principal then $\chi_1^s=\chi_2^\sigma$,
which is excluded. Retaining at $T(w)$ the zero of $\chi_i^{-\operatorname{sgn}w_i}$, the principal term at $w$
and the extra retained term together change the inequality by at most
$a_w\{\Re F(-a+iy)-\frac16f(0)-4t\Re F(\lambda_i-a+iy)\}\le0$, by \eqref{eq:alias}. (ii) Use
$P_{t'}(\chi_1(n)n^{-i\gamma_1})\bigl(1+\Re\chi_2(n)n^{-i\gamma_2}\bigr)\ge0$. Now $\chi_2$ is real, with coefficient
$1$ and $\varphi=\frac14$; the only possible principal indices are $(\pm2,\pm1)$, which are matched to
$(\pm1,\pm1)$ with target $\overline{\rho_1}$ or $\rho_1$, as before.
\end{proof}

\begin{proposition}[Additional-zero bounds for a nonreal first character]\label{prop:polyrows}
Let $\chi_1$ be nonreal of order at least $5$, let $\lambda_1\in[a,b]$ with $0\le a\le b$, and suppose
$\lambda'\le h$. Fix $t,\gamma>0$, let $c_1,c_2$ be the coefficients of $P_t$, and let
$F$ be the transform of $f=f_\gamma$. Choose an admissible $\rho'$ of $\chi_1$ with
parameter $\lambda'$ and height $\gamma'$. Put
$x=\lambda_1-a$, $y=\lambda'-a$ and $\Delta=\LL(\gamma_1-\gamma')$, and
\[
  M_t=(1+c_1+c_2)^2-1-\tfrac12(c_1^2+c_2^2),\quad p_1=\tfrac12c_1^2,\quad p_2=\tfrac12c_2^2,\quad
  q_1=c_1(1+\tfrac12c_2),\quad q_2=\tfrac12c_1c_2 .
\]
Then for every $\eps>0$ and all sufficiently large $q$,
\[
  c_1[F(x)+F(y)]\le F(-a)+\tfrac{M_t}6f(0)+R(\Delta)+\eps,
\]
where
\begin{align*}
 R(\Delta)={}&p_1\Re F(-a+i\Delta)+p_2\Re F(-a+2i\Delta)\\
 &-q_1\Re\{F(x+i\Delta)+F(y+i\Delta)\}\\
 &-q_2\Re\{F(x+2i\Delta)+F(y+2i\Delta)\}.
\end{align*}
\end{proposition}

\begin{proof}
Apply \cref{lem:positivity} to $P_t(\chi_1(n)n^{-i\gamma_1})P_t(\chi_1(n)n^{-i\gamma'})\ge0$, whose terms are
$\chi_1^{k+l}$ at heights $k\gamma_1+l\gamma'$, $|k|,|l|\le2$. Since the order of $\chi_1$ is at least $5$ and
$|k+l|\le4$, the principal terms are exactly those with $k+l=0$; they have total mass $p_1$ at height
$\pm\Delta$ and $p_2$ at $\pm2\Delta$. In every term with $k+l=\pm1$ retain both $\rho_1$ and $\rho'$ (or their
conjugates); a multiple zero is retained with its multiplicity. The terms $(1,0),(0,1)$ and their
conjugates give $-c_1[F(x)+F(y)]-c_1\Re[F(x+i\Delta)+F(y+i\Delta)]$, and the terms $(2,-1),(-1,2)$ and their
conjugates give $-\frac12c_1c_2\Re[F(x+i\Delta)+F(y+i\Delta)+F(x+2i\Delta)+F(y+2i\Delta)]$.
\end{proof}

For $\chi_1$ of order $3$ or $4$, \cref{inp:complextables}(b) gives $\lambda'>1.35$ whenever
$\lambda_1\le0.86$, which is stronger than every bound in \cref{prop:polyrows} that we use.

The rows used are the following, each on a cell of width $0.0025$ and each with a positive
normalized margin. There are $25$ rows of \cref{prop:complexrows}, on cells covering
$[0.66,0.7225]$, with conclusions from $\lambda_2>0.763$ (on $[0.66,0.6625]$) to $\lambda_2>0.725$ (on
$[0.72,0.7225]$), parameters $\gamma\in\{1.25,1.3\}$, $t\in\{0.9,0.95\}$ and, for a real $\chi_2$,
test parameter $\gamma_{\mathrm r}\in\{1.1,1.15\}$ and $t'=0.85$; their generic margins are at least $8.2\cdot10^{-4}$, and the margins
in \eqref{eq:alias} at least $0.163$. There are $69$ rows of \cref{prop:polyrows}, on $69$ of the $70$ cells of width $0.0025$ in
$[0.68,0.855]$ (all except $[0.6975,0.70]$), with conclusions from $\lambda'>0.975$ (on $[0.68,0.6825]$) to $\lambda'>0.855$ (on
$[0.8525,0.855]$), parameters $\gamma\in\{1.05,1.1,1.15\}$ and $t\in\{0.9,0.95\}$, and margins at least
$1.5\cdot10^{-3}$. For these rows, if $\lambda'\le h$, the left side of \cref{prop:polyrows} is at least
$c_1[F(b-a)+F(h-a)]$, since $F$ is decreasing, and $R(\Delta)$ is replaced by its supremum over all real
$\Delta$, $x\in[0,b-a]$ and $y\in[p-a,h-a]$, where $p=\max(p_j,a)\le\lambda'$ comes from the bound $p_j$ of the
parent row of the cell (\cref{prop:parentrows}; it comes from \cite[Table~$2'$]{Xylouris2011nullstellen}) and the
trivial bound $\lambda'\ge\lambda_1\ge a$: $p=0.96$, $0.93$, $0.91$, $0.89$, $0.86$, $0.84$, $0.83$ on the cells in $[0.68,0.70]$,
$[0.70,0.72]$, \dots, $[0.80,0.82]$, $p=0.827$ on $[0.82,0.8275]$, and $p=a$ on the cells beyond. So each row proves
that $\lambda_1\in[a,b]$ and $\lambda'\ge p$ imply $\lambda'>h$, and since $\lambda'\ge p$ holds for every configuration
with $\lambda_1\in[a,b]$, $\lambda_1\in[a,b]$ implies $\lambda'>h$. The complete list is in Appendix~\ref{app:rows}. On $11$ further cells the inequalities of
\cref{prop:complexrows} are evaluated on the actual cell of the case tree.

\subsection{A positivity exclusion}

\begin{proposition}[Excluding a nearby real second character]\label{prop:positivity}
Let $\chi_1,\rho_1$ be real with $\lambda_1\in[a,b]$, and let $\chi_2\ne\chi_1$ be a real nonprincipal
character with a zero of height at most $1$ and parameter $\nu\le h$. If
\[
  F(b-a)+F(h-a)-F(-a)-\tfrac38f(0)>0
\]
for a test function \eqref{eq:ftest}, then this situation does not occur for $q\ge q_0$.
\end{proposition}

\begin{proof}
Apply \cref{lem:positivity} to $(1+\chi_1(n))(1+\Re\chi_2(n)n^{-i\gamma_2})\ge0$. Its three nonprincipal
characters $\chi_1$, $\chi_2$, $\chi_1\chi_2$ are real, so $\varphi=\frac14$; retain $\rho_1$ and the given zero of
$\chi_2$, which lies in $R(l)$. This gives $F(\lambda_1-a)+F(\nu-a)\le F(-a)+\frac38f(0)+\eps$.
\end{proof}

It excludes $197$ nodes of the case tree (\cref{subsec:vocabulary}), all of type $\mathrm{rr}$ with
a reserved real family; the smallest normalized margin is $7.8\cdot10^{-5}$.

\subsection{Applying the rows}

\begin{lemma}[Preservation of configurations under location updates]\label{lem:updates}
In the notation of \cref{def:specification}, suppose that a result of this section proves $\lambda_2>h$ for every configuration of type $\mathfrak t$ with
$\lambda_1\in[a',b']$, and let $\mathfrak s$ be a specification of type $\mathfrak t$ with cell $[a,b]\subseteq[a',b']$.
Then replacing $l_2$ and $r$ by $\max(l_2,h)$ and $\max(r,h)$ (and $\lo_2$ by the new $r$) does not
change $C(\mathfrak s)$, and $C(\mathfrak s)=\emptyset$ if $\mathfrak s$ is reserved with $\hi_2\le h$. Similarly a proof of
$\lambda'>h$ allows $p$ and $\lo'$ to be replaced by $\max(\cdot,h)$, and makes $C(\mathfrak s)$ empty if
$\hi'\le h$.
\end{lemma}

\begin{proof}
Every family other than the first has all its zeros in $R(l)$ at parameter at least $\lambda_2>h$; in
particular $\nu(\mathcal F)>h$.
\end{proof}

\subsection{The third family}

\begin{proposition}[Lower bounds for the third family]\label{prop:third}
Consider a configuration satisfying a specification of type $\mathfrak t$ with
$\lambda_1\in[a,b]$. Put $h=\hi_2$ if a family is reserved and $h=\infty$ otherwise;
then $\lambda_2\le h$ by \cref{lem:familyfacts}.
Each applicable rule below gives a strict lower bound for $\lambda_3$, for all
sufficiently large $q$. Let $\lambda_3^{\lo}$ be the maximum of the finitely many bounds
selected for the leaf. Then $\lambda_3>\lambda_3^{\lo}$.
\begin{enumerate}[(1)]
  \item $0.857$ (\cref{inp:lambda3}(a));
  \item for $\mathfrak t=\mathrm{rr}$ and $[a,b]$ inside $[0.44,0.60]$, $[0.60,0.70]$ or $[0.70,0.80]$: $1.176$,
    $1.055$, $0.952$ (\cref{inp:lambda3}(c));
  \item for $\mathfrak t\ne\mathrm{rr}$ and $b\le t$: the value $c$ of the smallest applicable row of
    \cref{inp:lambda3}(b);
  \item for $\mathfrak t=\mathrm{complex}$ and $0.44\le a\le b\le0.85$: any $t$ with
    \[
      F(t-b)+F(\min(t,h)-a)-F(-b)+F(0)-\tfrac76f(0)>10^{-5}f(0)\qquad(\gamma=\tfrac54);
    \]
  \item for $\mathfrak t=\mathrm{rr}$, $0.44\le a\le b\le0.80$ and $t\in[0.952,1.176]$: any $t$ such that
    $[a,\min(t,h)]$ is covered by intervals $[c,c']$ with
    \[
      F(-c')-F(t-c')-F(0)-F(b-c)+\tfrac98f(0)<-10^{-5}f(0)\qquad(\gamma=\tfrac{26}{25}).
    \]
\end{enumerate}
\end{proposition}

\begin{proof}
(1)--(3) are printed. For (4), suppose $\lambda_3\le t$; then $\lambda_2\le\min(t,h)$. The function
$\lambda\mapsto F(-\lambda)-F(t-\lambda)=\int f(u)e^{\lambda u}(1-e^{-tu})\,du$ is increasing, so
$F(-\lambda_1)-F(\lambda_3-\lambda_1)\le F(-\lambda_1)-F(t-\lambda_1)\le F(-b)-F(t-b)$, and
$F(\lambda_2-\lambda_1)\ge F(\min(t,h)-a)$. So the right side of \cref{inp:lambda3}(d) is below
$-10^{-5}f(0)+\eps<0$, a contradiction; the side condition \cite[(4.29)]{Xylouris2011nullstellen} holds by our verification. For (5),
suppose $\lambda_3\le t\le1.176$. Then $\lambda_2\in[a,\min(t,h)]\subseteq[0.44,1.176]$ and
$\lambda_1\in[0.44,0.80]$, so \cref{inp:lambda3}(e) applies. If $\lambda_2\in[c,c']$, the same monotonicity gives
$F(-\lambda_2)-F(\lambda_3-\lambda_2)\le F(-c')-F(t-c')$ and $F(\lambda_1-\lambda_2)\ge F(b-c)$, and the right side of
(e) is negative for small $\eps$, a contradiction; the side condition \cite[(4.34)]{Xylouris2011nullstellen} holds
by our verification.
\end{proof}

For example, on the $\mathrm{rr}$ cell $[0.7025,0.705]$ rule (5) gives $\lambda_3>1.0502$ without a
reservation, and $\lambda_3>1.0888$ with a reserved family below $0.85$; Table~10 gives $0.952$.
Rules (4) and (5) are evaluated with interval arithmetic, and a value $t$ is accepted only with a
normalized margin exceeding $10^{-5}$.
